\section{Conclusion}
We presented \model, a plug-in framework for sequential recommendation that
combines multi-view text features, explicit cross interactions, and contrastive
alignment. Our full-ranking experiments yield actionable practitioner guidance:

\textbf{(1) TF-IDF is surprisingly strong---start there.}
\tfidf alone provides substantial gains over ID-only baselines on both
overall and ranking metrics. This lightweight baseline should be the
\emph{first} text enhancement before investing in heavier \llm embeddings.

\textbf{(2) Multi-view benefits cold-start items with LLM scale effects.}
\model achieves best overall performance on both datasets, with the most
pronounced gains on the \textbf{new} stratum (cold-start items). Furthermore,
larger LLM encoders show clear scaling behavior on long-tail items: as model
size increases, cold-start coverage improves consistently. When cold-start
performance is critical, multi-view prompting with larger LLMs justifies
the embedding cost; for simpler catalogs, \tfidf alone may suffice.

\textbf{(3) Cross network controls the recall--precision trade-off.}
Toggle cross based on deployment stage:
\emph{candidate generation} (maximize Top-$K$ coverage)---disable cross for
higher HR at the cost of lower MRR;
\emph{re-ranking} (maximize Top-1 quality)---enable cross to sharpen score
distributions and improve MRR/NDCG.

\textbf{(4) Tune \texttt{infer\_boost} at deployment---no retraining needed.}
Sensitivity analysis shows \texttt{infer\_boost} is the most impactful
hyperparameter for cold-start performance, while contrastive hyperparameters
($\lambda$, $\tau$) are robust across reasonable ranges.
\textbf{Recommendation}: use moderate \texttt{infer\_boost} for balanced
performance; increase it when cold-start items are prioritized.

\bibliographystyle{unsrt}
\bibliography{refs}

\clearpage  % Force new page for appendix

\appendix
